\chapter*{I. Introducción}

¿Cómo afectan las temperaturas extremas a los grandes del café... y a los más pequeños? 

Sin duda, el cambio climático representa una amenaza crítica para la producción agrícola global, con grave impacto sobre los \textit{soft commodities} --o productos básicos de cultivo--, que necesitan condiciones estables de temperatura y precipitaciones para su correcto desarrollo. Esto es particularmente importante si se considera que la agricultura continúa siendo una de las principales fuentes de ingresos para países en desarrollo [CITAR]

En ese sentido, el pronóstico es desafiante, toda vez que, producto del cambio climático, se espera que aumente notoriamente la frecuencia de eventos meteorológicos extremos, llevando a un aumento en la temperatura global y a alteraciones en los patrones de precipitaciones [CITAR]. Dado que los rendimientos de los cultivos son susceptible a estos cambios, es posible que se vean perjudicados cuanto mayor sea la transformación de los factores climáticos.

Los cafetos, por ejemplo, prosperan bajo condiciones ambientales específicas: son sensibles a cambios en los patrones de lluvia y pueden verse gravemente perjudicados por la exposición prolongada a temperaturas extremas [CITAR]. En consecuencia, grandes desviaciones en la trayectoria histórica de estas series pueden tener efectos negativos sobre el desarrollo del grano de café, lo que se verá reflejado en el precio y en las cantidades comercializadas en el mercado internacional [CITAR].

En las últimas décadas, la estimación rigurosa del impacto del cambio climático sobre diversas variables económicas se ha consolidado como una línea de investigación de gran relevancia en la literatura económica. El desarrollo metodológico para cuantificar estos efectos ha permitido avanzar en la comprensión de las implicaciones económicas del fenómeno climático, distinguiendo entre efectos de corto, mediano y largo plazo, e identificando medidas de mitigación por parte de los agentes afectados.

Por un lado, el uso de modelos de datos de panel permite explotar las fluctuaciones interanuales o de corto plazo, aislando el efecto del fenómeno climático de las características fijas de cada país [CITAR]. Así, por ejemplo, se ha estimado que un aumento de 1°C en la temperatura local puede reducir el crecimiento económico anual en países pobres aproximadamente 1,3 puntos porcentuales [CITAR].

Por otro lado, en contraste con un análisis de choques de corto plazo, otra línea de investigación se ha centrado en evaluar la capacidad de adaptación y mitigación de diferentes agentes en el mediano y largo plazo. El método de largas diferencias (\textit{Long Differences}) compara promedios de variables en dos puntos distantes en el tiempo, y ha sido empleado para estudiar cómo los agentes responden a cambios graduales o permanentes en el clima. Este enfoque, por ejemplo, ha permitido identificar que las adaptaciones de largo plazo en China mitigaron los impactos negativos del calor extremo en la productividad total de factores en un 37,9\%, y en el rendimiento agrícola en un 46,8\% [CITAR].

Otra distinción importante en la estimación de impacto es la identificación del mecanismo a revisar, ya que el fenómeno climático y sus consecuencias pueden abarcar diferentes elementos en la cadena de producción de los \textit{soft commodities} y en sus respectivos mercados. En tal sentido, no solo la elección de la metodología es relevante; la contribución de cada investigación también puede cambiar según las variables de interés involucradas en el análisis.

La presente investigación tiene como objeto de análisis el caso específico del mercado internacional del café, y los efectos de la exposición prolongada a temperaturas extremas (o estrés térmico). En particular, se tienen cuatro objetivos específicos. Primero, estimar estructuralmente el mercado internacional del café, mediante el uso de variables instrumentales para la demanda global y un modelo simple de costos marginales de producción. Segundo, simular escenarios contrafactuales sin estrés térmico y cuantificar su impacto sobre las variables del mercado y el bienestar social. Tercero, evaluar diferencias de impacto que pudieran existir entre países líderes y seguidores, distinguiendo entre las que son inherentes a la estructura del mercado y aquellas provocadas por el estrés térmico.

[RESUMEN DE RESULTADOS]

[RESUMEN DE SECCIONES]